% Section 9. Input from paper.tex; every number traces to sources.md and is asserted by check_numbers.py.
\providecommand{\qs}{QScore9}

\section{Registration and reproducibility}
\label{sec:repro}

The repository \url{https://github.com/vprchlik/PulseGraph} holds the protocol with its dated additions,
whose registered bytes are byte prefixes of the file as committed and were verified by each run's gate; the
declaration files of every family; the study code with its offline tests; and, for each study, the evidence
package: the frozen frame, manifest and selection, the panel's status file, the readiness and gate records,
the run's state with the SHA-256 of every checkpoint, the inference directory with the family table, the
per-series paired differences and the secondary tables, the integrity recomputation, the results page and a
provenance file with the approvals, the commands and their times. Every score of the three runs was
recomputed from the saved forecasts by a separate script with no disagreement, and the Study 3 census was
recomputed from the raw captures by a script that shares no code with the pipeline. Each post-hoc analysis
has its own directory with its specification, command, outputs and provenance. Not in the tree are the
checkpoints that hold every forecast, 300 for GitHub (one results file of 433,172,286 bytes), 540 for the
grid, 209 for npm and 300 for NYC 311, the raw captures, the target signals and the daily counts of the npm
and NYC 311 series; the checkpoints and the aggregated daily counts are to be released as release assets
beside the paper, under the data statements of Appendices~\ref{app:s2panel} and~\ref{app:s3detail}, so that
a fresh clone can rerun every inference, integrity check and post-hoc analysis. The commands are given by
reference to each run's results page in Appendix~\ref{app:runs}, with the run record; the number checker
asserts every figure in this paper against the files it names, and the figure script draws every figure
from the saved files alone.
